\clearpage
\begin{trapbox}
    
\begin{center}\begin{tikzpicture}[font=\small,>=stealth,line join=round]
\draw[gray!45,rounded corners=3pt] (0,0) rectangle (14.2,9.4);
\node at (7.1,8.8) {单调但不连续，端点之间可能存在整段缺值};

\begin{scope}[shift={(2.5,2.3)},x=4cm,y=2.3cm]
\draw[->] (-.1,0)--(1.4,0) node[right]{$x$};\draw[->] (0,-.1)--(0,2.5) node[above]{$y$};\fill[orange!12] (0,1) rectangle(1,2);
\draw[blue,very thick] (0,0)--(1,1);\fill[blue] (1,2) circle(2.5pt);\fill[blue] (0,0) circle(2pt);
\draw[dashed] (0,2)--(1,2)--(1,0);\node[below] at (1,0){$1$};\node[left] at (0,1){$1$};\node[left] at (0,2){$2$};
\draw[blue,fill=white,thick] (1,1) circle(2.5pt);
\end{scope}
\node[align=left,anchor=west] at (8.7,5.1){$f(x)=x\quad(0\le x<1)$\\[10pt] $f(1)=2$\\[12pt] 最小值 $0$，最大值 $2$\\[10pt] 值域 $[0,1)\cup\{2\}$};
\node[text width=13.2cm,align=center] at (7.1,.65) {浅橙色区域对应未取到的高度；$y=1$ 也没有取到。};
\end{tikzpicture}
\end{center}
\begin{description}[style=nextline,leftmargin=2.8em,labelsep=0.5em,itemsep=0.48em]
        \item[\textbf{\textcolor{CTrap}{易错一：只凭单调，就写出完整的闭区间值域}}]
            \textbf{反例：}设函数在 $[0,1]$ 上满足
            \[
                f(x)=\begin{cases}
                    x,&0\le x<1,\\
                    2,&x=1.
                \end{cases}
            \]
            它确实一直递增，$f(0)=0$、$f(1)=2$，最小值和最大值都在端点。
            可是真正的值域是
            \[
                \boxed{[0,1)\cup\{2\}},
            \]
            \textbf{不是} $[0,2]$。因为 $x$ 走到 $1$ 时，$y$ 从接近 $1$ 突然跳到 $2$，
            中间 $1$ 到 $2$ 之间的数都没取到。
            \par\textbf{正确做法：}完整值域公式还需保证整段图像连续。

        \item[\textbf{\textcolor{CTrap}{易错二：开区间也把端点值算进去}}]
            \textbf{反例：}$f(x)=x$，$x\in(0,2)$。
            这个函数递增且连续，但 $x$ 不能取 $0$ 或 $2$，所以值域是 $(0,2)$，而不是 $[0,2]$。
            \par\textbf{正确做法：}先检查两端能不能取到。不能取到的端点不可直接写方括号。

        \item[\textbf{\textcolor{CTrap}{易错三：函数中途转弯，仍然只代两端}}]
            \textbf{反例：}$f(x)=x^2$，$x\in[-1,2]$。
            $f(-1)=1$、$f(2)=4$，但 $f(0)=0$ 更小，
            所以值域是 $[0,4]$，不是 $[1,4]$。
            \par\textbf{正确做法：}先看是否跨过二次函数的对称轴，必要时分段。

        \item[\textbf{\textcolor{CTrap}{易错四：递减时把左右顺序写反}}]
            \textbf{反例：}$f(x)=3-x$，$x\in[1,4]$。
            左端函数值是 $2$，右端函数值是 $-1$；值域应写 $[-1,2]$。
            \par\textbf{正确做法：}无论递增、递减，区间都要写\textbf{小的在左，大的在右}。

        \item[\textbf{\textcolor{CTrap}{易错五：定义域中存在不能代入的点}}]
            \textbf{反例：}$f(x)=\dfrac{1}{x}$，题目要求研究 $[-1,1]$。
            因为 $x=0$ 时不能计算，这个函数根本没有在整个 $[-1,1]$ 上定义，
            更谈不上在这一整段上直接使用端点法。
            \par\textbf{正确做法：}先检查分母能不能为零、根号内是否符合要求等，
            再分别考虑真正可以取值的区间。
    \end{description}
\end{trapbox}
